\section{Sharpness of \texorpdfstring{$\hG^{3/2}$}{hG^{3/2}}}\label{sec:sharp}

Theorem~\ref{er:D} gives $\norm{\nabla(\ut-u_h)}\le C[(1+\gamma^{1/2})\hG^{3/2}+\epsd]$ for $\CNS$, and Theorem~\ref{er:A} the same for $\GS$ when $G=0$. We prove the reverse inequality whenever $|\hh|\,\partial_Nu\not\equiv0$ on $\Gamma$. The lower bound does not use (IS3), except to guarantee that $u_h$ exists.

The idea is to test the error equation with discrete divergence-free fields $v$ that vanish on $\Gh$. For such $v$ neither the pressure nor the penalty enters, and the error equation says that $\ut-u_h$ must compensate the Nitsche term $\ip{\ut}{\dn v}$, which is nonzero because $\ut$ does not vanish on $\Gh$. On a face, $\ut\approx d\,\partial_Nu$ with $d$ the quadratic sagitta, so this term is a sum of three edge moments of $\dn v$ per face, weighted by the second fundamental form. The lower bound follows once we have local fields, one per face, that make these moments as large as scaling allows. This is the local condition (M3$^3$) below; the rest of the section verifies it.

\subsection{An abstract lower bound}
Let $Y_h:=\Zh\cap\VO$ and
\[
Y_h^1:=\Bigl\{v\in Y_h:\ \int_F\dn v\,dA=0\ \ \forall F\in\FG\Bigr\}.
\]
For $F\in\FG$ let $T_F$ be the tetrahedron of $\Th$ with face $F$; the $T_F$ are distinct (Lemma~\ref{geom:one}).

\begin{lemma}\label{sh:normal}
If $v\in Y_h$ and $F\in\FG$, then on $F$ (trace from $T_F$) $\nabla v=\dn v\otimes n_F$ and $\dn v\cdot n_F=0$.
\end{lemma}

\begin{proof}
$v|_{T_F}$ is a polynomial vanishing on the plane of $F$, so its tangential derivatives vanish there, and $0=\div v=\dn v\cdot n_F$.
\end{proof}

\begin{lemma}[Error identity on $Y_h$]\label{sh:identity}
Let $u_h$ be the velocity of $\CNS$ ($\gamma\ge\gamma_2$) or of $\GS$ ($\gamma\ge\gamma_0$, any $G$), and $e_u:=\ut-u_h$. For $v\in Y_h$, $a_h(e_u,v)-\ip{e_u}{\dn v}=-\ip{\ut}{\dn v}-(F,v)$.
\end{lemma}

\begin{proof}
By \eqref{er:cnserr}, resp.\ Corollary~\ref{er:GSZ}, $N_h(e_u,v)=\Gc(v)$ on $Y_h$: the pressure terms vanish because $\div v=0$ and $v=0$ on $\Gh$. Since $v=0$ on $\Gh$, $N_h(e_u,v)=a_h(e_u,v)-\ip{e_u}{\dn v}$, and the transposed and penalty parts of $\Gc(v)$ vanish.
\end{proof}

\begin{lemma}[Trace control on $Y_h^1$]\label{sh:poinc}
Let $C_{PT}$ be the constant in the scaled Poincar\'e trace inequality $\norm{w-\bar w_T}_{L^2(F)}\le C_{PT}(\diam T)^{1/2}\norm{\nabla w}_T$ for a face $F$ of a tetrahedron $T$, where $\bar w_T$ is the mean of $w$ over $T$. Then $|\ip w{\dn v}|\le C_{PT}C_I\norm{\nabla w}\norm{\nabla v}$ for $w\in H^1(\Oh)^3$ and $v\in Y_h^1$.
\end{lemma}

\begin{proof}
$\int_Fw\cdot\dn v=\int_F(w-\bar w_{T_F})\cdot\dn v$; use Cauchy--Schwarz on each $F$, the local form of Lemma~\ref{st:inverse}, and Cauchy--Schwarz over the distinct $T_F$.
\end{proof}

\begin{corollary}[Abstract lower bound]\label{sh:abstract}
With $\Lambda_0:=2+C_{PT}C_I$,
\[
\Lambda_0\norm{\nabla e_u}\ge\sup_{0\ne v\in Y_h^1}\frac{|\ip\ut{\dn v}+(F,v)|}{\norm{\nabla v}} .
\]
\end{corollary}

\begin{proof}
$|a_h(e_u,v)|\le2\norm{\nabla e_u}\norm{\nabla v}$, and Lemmas~\ref{sh:identity} and~\ref{sh:poinc}.
\end{proof}

Neither the pressure nor the penalty enters: the identity is the discrete momentum equation on $Y_h$, which $\ut$ violates through Nitsche's symmetric term acting on the nonzero trace of $\ut$ on $\Gh$.

\subsection{The leading term: three edge moments per face}
Let $W:=\partial_Nu|_\Gamma$, a tangential field (Lemma~\ref{geom:wall}), and $W_F:=W(y_F)$. For $F\in\FG$ put $w^F_{ij}:=\hh_F(e_{ij},e_{ij})$, $Q_F:=\frac12\sum_{i<j}w^F_{ij}\mu_i\mu_j$ (so $d=Q_F+O(\ell^3)$ on $F$ by Lemma~\ref{geom:face}(a)), and for $v\in Y_h$
\[
\M_F(v):=\int_FQ_F\,\dn v\,dA=\frac12\sum_{i<j}w^F_{ij}\int_F\mu_i\mu_j\,\dn v\,dA,\qquad \ell_W(v):=\sum_{F\in\FG}W_F\cdot\M_F(v).
\]
By Lemma~\ref{sh:normal}, $\M_F(v)$ is parallel to $F$. In the model case $w^F_{ij}=-\ell_{ij}^2$ (Lemma~\ref{geom:sag}).

\begin{lemma}[Leading term]\label{sh:lead}
For $v\in Y_h$, $|\ip\ut{\dn v}-\ell_W(v)|\le C\hG^{5/2}\norm{\nabla v}$ and $|(F,v)|\le C\hG^2\norm{\nabla v}$.
\end{lemma}

\begin{proof}
On $\Gh$, $x=x^\ast+d(x)N(x^\ast)$ and $\ut(x^\ast)=0$, so $\ut(x)=d(x)W(x^\ast)+O(d^2)$. With $|d-Q_F|\le C\hG^3$, $|W(x^\ast)-W_F|\le C\hG$ and $|Q_F|\le C\hG^2$, $\ut=Q_FW_F+R$ on $F$ with $|R|\le C\hG^3$. Hence $\int_F\ut\cdot\dn v=W_F\cdot\M_F(v)+O(\hG^3)\norm{\dn v}_{L^1(F)}$, and $\sum_F\norm{\dn v}_{L^1(F)}\le\sum_F|F|^{1/2}C_I(\diam F)^{-1/2}\norm{\nabla v}_{T_F}\le C\hG^{-1/2}\norm{\nabla v}$ since $\#\FG\le C\hG^{-2}$. The bound on $(F,v)$ is Lemma~\ref{geom:ext}.
\end{proof}

\subsection{The local condition and the lower bound}
\begin{assumption}[(M3$^3$)]\label{sh:M3}
There are $\kappa_0>0$ and $K_{\rm ov},C_s,C_1$ such that for every $F\in\FG$ and every unit $t\in T_{y_F}\Gamma$ there is $v_{F,t}\in Y_h^1$ with
\begin{enumerate}[label=(\roman*)]
\item $\norm{\nabla v_{F,t}}=1$;
\item $\operatorname{supp}v_{F,t}\subset\omega_F$, a union of tetrahedra within distance $C_s\hG$ of $F$, each tetrahedron lying in at most $K_{\rm ov}$ of the sets $\omega_F$;
\item $t\cdot\sum_{F'\subset\omega_F\cap\Gh}\M_{F'}(v_{F,t})\ \ge\ \kappa_0\,|\hh_{y_F}|\,(\diam F)^{5/2}-C_1(\diam F)^{7/2}$.
\end{enumerate}
\end{assumption}
Condition (iii) is scale invariant: under $x\mapsto\lambda x$ (with $v$ renormalised) $\norm{\nabla v}$ scales like $\lambda^{1/2}$, $w^F$ like $\lambda^2$ and $\M_F$ like $\lambda^3$. Since (iii) is linear in $v$, only the sign of $v$ is at our disposal; statements about $-\M_F$ are equivalent. The vector $t$ is tangent to $\Gamma$ at $y_F$, not to $F$; this is the form in which the witnesses are assembled.

Given the local witnesses, the lower bound is obtained by adding them up, face by face, with weights $|\hh|\,|W|$ and directions $W/|W|$; the overlap of their supports is bounded, so their energies add.

\begin{theorem}[$H^1$ lower bound]\label{sh:lb}
Assume the setting of Section~\ref{sec:setting}, (M3$^3$), and $|\hh|\,W\not\equiv0$. Let $u_h$ be the velocity of $\CNS$ ($\gamma\ge\gamma_2$) or of $\GS$ ($\gamma\ge\gamma_0$, any $G$). Then for $\hG\le h_1$
\[
\norm{\nabla(\ut-u_h)}_{\Oh}\ \ge\ c\,\hG^{3/2}\,\bigl\|\,|\hh|\,W\bigr\|_{L^2(\Gamma)}-C\hG^2,
\]
with $c=c(\kappa_0,K_{\rm ov},\Lambda_0,c_0,\theta_0)$ and $C$ independent of $\gamma,\mu,\rho,h$, the pressure and $\epsd$. Here $h_1$ depends also on $C_s$, $C_1$, $\norm W_{W^{1,\infty}}$ and $\norm{|\hh|W}_{L^2}$. In the model case $|\hh|=1$.
\end{theorem}

\begin{proof}
For each $F$ with $W_F\ne0$ let $t_F:=W_F/|W_F|\in T_{y_F}\Gamma$ and $v_F:=|\hh_{y_F}|\,|W_F|\,v_{F,t_F}$ (else $v_F:=0$), and $v:=\sum_Fv_F\in Y_h^1$. By linearity,
\[
\ell_W(v)=\sum_F\Bigl[W_F\cdot\sum_{F'}\M_{F'}(v_F)+\sum_{F'}(W_{F'}-W_F)\cdot\M_{F'}(v_F)\Bigr],
\]
the inner sums running over the (boundedly many) faces $F'\subset\omega_F$. Since $|\M_{F'}(v)|\le C\hG^2|F'|^{1/2}\norm{\dn v}_{L^2(F')}\le C\hG^{5/2}\norm{\nabla v}$ and $|W_{F'}-W_F|\le CC_s\hG\norm W_{W^{1,\infty}}$, (M3$^3$)(iii) gives
\[
\ell_W(v)\ge\kappa_0c_0^{5/2}\hG^{5/2}\sum_F|\hh_{y_F}|^2|W_F|^2-C\hG^{7/2}\sum_F|\hh_{y_F}||W_F|\bigl(|W_F|+\norm W_{W^{1,\infty}}\bigr),
\]
and $\norm{\nabla v}^2\le K_{\rm ov}\sum_F|\hh_{y_F}|^2|W_F|^2$. The projections $\pi(F)$ tile $\Gamma$ with area Jacobian $1+O(\hG^2)$ (Theorem~\ref{geom:orient}, Lemma~\ref{geom:face}(c)), and $|\hh|,|W|$ are Lipschitz, so $S:=\sum_F|F||\hh_{y_F}|^2|W_F|^2=\norm{|\hh|W}^2_{L^2(\Gamma)}+O(\hG)$; with $c_A\hG^2\le|F|\le\hG^2$ ($c_A$ from $c_0,\theta_0$) we get $\hG^{-2}S\le\sum_F|\hh_{y_F}|^2|W_F|^2\le(c_A\hG^2)^{-1}S$, and the error sum is at most $C\#\FG\le C\hG^{-2}$. Hence
\[
\frac{\ell_W(v)}{\norm{\nabla v}}\ge\kappa_0c_0^{5/2}\Bigl(\frac{c_A}{K_{\rm ov}}\Bigr)^{1/2}\hG^{3/2}S^{1/2}-C\frac{\hG^{5/2}}{S^{1/2}},
\]
and for $\hG\le h_1$ the second term is at most half the first and $S\ge\frac34\norm{|\hh|W}^2$. Conclude with Corollary~\ref{sh:abstract} and Lemma~\ref{sh:lead}.
\end{proof}

On flat parts of $\Gamma$ the faces lie in $\Gamma$ and the chordal bump vanishes, so the weight $|\hh|$ cannot be removed.

\subsection{Witnesses by affine transport}\label{sh:sec:transport}
On a single macro-element, witnesses can be computed once on a reference tetrahedron and transported to every shape by the Piola map, which preserves divergence-free fields and transforms the edge moments by an explicit linear map. So (M3$^3$) reduces to a finite-dimensional surjectivity statement on the reference element, which can be checked in exact arithmetic.

Let $T$ be a tetrahedron with a face $F$ and apex $p_0$, $\hat T=\conv(\hat p_0,\dots,\hat p_3)$ with $\hat p_0=e_3$, $\hat p_1=0$, $\hat p_2=e_1$, $\hat p_3=e_2$ and $\hat F=\hat T\cap\{\hat x_3=0\}$, and $x=A\hat x+a$ the affine map with $\hat F\mapsto F$, $\det A>0$. Let $\nu=-n_F$ be the inward unit normal of $T$ on $F$ and $h_T$ the height of $T$ over $F$. For $\hat v$ vanishing on $\hat F$ let $v:=(\det A)^{-1}A\,\hat v\circ A^{-1}$ (Piola) and $\gh:=\partial_{\hat\nu}\hat v|_{\hat F}$.

\begin{lemma}[Piola transport]\label{sh:piola}
(a) The Piola map preserves continuity, zero boundary values and piecewise-$P_k$ structure on affinely corresponding subdivisions, and $\div v=(\det A)^{-1}(\widehat\div\,\hat v)\circ A^{-1}$; it maps the Alfeld split of $\hat T$ onto that of $T$.
(b) For every polynomial weight $\phi$ in the barycentrics of $F$, $\int_F\phi\,\partial_\nu v\,dA=h_T^{-2}A\int_{\hat F}\phi\,\gh\,d\hat A$.
(c) $\norm{\nabla v}_T\le(\det A)^{-1/2}\norm A\norm{A^{-1}}\norm{\hat\nabla\hat v}_{\hat T}$.
\end{lemma}

\begin{proof}
(a) is standard ($A$ is constant and maps barycentre to barycentre). (b) On $\hat F$, $\hat\nabla\hat v=\gh\otimes e_3$, so $\partial_\nu v=(\det A)^{-1}A\gh\,(A^{-\mathsf T}e_3\cdot\nu)$ with $A^{-\mathsf T}e_3=|A^{-\mathsf T}e_3|\nu$. By Nanson's formula $dA=\det A\,|A^{-\mathsf T}e_3|\,d\hat A$, so the factor is $|A^{-\mathsf T}e_3|^2=|\nabla\lambda_0|^2=h_T^{-2}$. (c) $\nabla v=(\det A)^{-1}A(\hat\nabla\hat v)A^{-1}$.
\end{proof}
Since $\dn=-\partial_\nu$ on $F$, $\M_F(v)=-\frac12\sum w_{ij}\int_F\mu_i\mu_j\partial_\nu v$; the sign is absorbed by $v\mapsto-v$.

\begin{proposition}[Joint criterion]\label{sh:joint}
Let $\hat{\mathcal Y}$ be a space of continuous piecewise-polynomial divergence-free fields on a subdivision of $\hat T$ vanishing on $\partial\hat T$, $\hat{\mathcal Y}^1:=\{\hat v\in\hat{\mathcal Y}:\int_{\hat F}\gh=0\}$, and suppose the joint edge-moment map
\[
\hat E:\hat{\mathcal Y}^1\to(\R^2)^3,\qquad\hat v\mapsto\Bigl(\int_{\hat F}\mu_1\mu_2\gh,\ \int_{\hat F}\mu_1\mu_3\gh,\ \int_{\hat F}\mu_2\mu_3\gh\Bigr)
\]
is onto, with right inverse $\hat R$ (Euclidean norm on $\R^6$ to $\norm{\hat\nabla\cdot}_{L^2(\hat T)}$). Then for every tetrahedron $T$ with face $F$, every real weight vector $w\ne0$ and every unit $t\parallel F$ there is $v$ in the Piola image, supported in $T$, with $\norm{\nabla v}=1$ and
\[
t\cdot\frac12\sum_{i<j}w_{ij}\int_F\mu_i\mu_j\,\partial_\nu v\ \ge\ \frac{|w|\,(\det A)^{1/2}}{2h_T^2\norm A\norm{A^{-1}}^2\norm{\hat R}}\ \ge\ c(\text{shape})\,\frac{|w|\,(\diam F)^{1/2}}{\norm{\hat R}} .
\]
\end{proposition}

\begin{proof}
$s:=A^{-1}t$ is parallel to $\hat F$. Take $\hat v:=\hat R\bigl((w_{ij}s/|w|^2)_{i<j}\bigr)$, so that $\sum_{i<j}w_{ij}\hat E_{ij}(\hat v)=s$ and $\norm{\hat\nabla\hat v}\le\norm{\hat R}\norm{A^{-1}}/|w|$. By Lemma~\ref{sh:piola}(b), $\frac12\sum w_{ij}\int_F\mu_i\mu_j\partial_\nu v=\frac12h_T^{-2}As=\frac12h_T^{-2}t$; normalise with (c). Shape regularity gives $\norm A\lesssim\diam T$, $\norm{A^{-1}}\lesssim(\diam T)^{-1}$, $\det A\gtrsim(\diam T)^3$ and $h_T\lesssim\diam T\lesssim\diam F$.
\end{proof}
No compactness over shapes is used: $\hat{\mathcal Y}^1$ and $\hat E$ are fixed, and the shape enters only through $A$ and $w$.

\begin{theorem}[(M3$^3$) on Alfeld refinements, $k\ge4$]\label{sh:alfeld}
Let $\Th$ be an Alfeld refinement and $k\ge4$. Then (M3$^3$) holds with $\omega_F=T_F$ (the macro-tetrahedron containing $F$), $K_{\rm ov}=1$ and $\kappa_0=c(\text{shape regularity},\theta_0)/\norm{\hat R}$. \status{computer-assisted, exact}
\end{theorem}

\begin{proof}
Let $\hat{\mathcal Y}_k$ be the continuous piecewise-$P_k$ divergence-free fields on the Alfeld split of $\hat T$ vanishing on $\partial\hat T$. An exact rational computation in the Bernstein--B\'ezier basis ($35$ interior domain points, $105$ unknowns, $80$ divergence equations) gives, for $k=4$: the divergence has rank $79$ (onto the mean-free piecewise $P_3$), $\dim\hat{\mathcal Y}_4=26$, and the eight functionals ($\int_{\hat F}\gh$ and the three edge moments, two tangential components each) have rank $8$ on $\hat{\mathcal Y}_4$. So $\hat E$ is onto on $\hat{\mathcal Y}^1_4$ ($\dim=24$), and $\hat{\mathcal Y}_4\subset\hat{\mathcal Y}_k$. By Lemma~\ref{sh:piola}(a),(b) the Piola image on $T_F$, extended by zero, lies in $Y_h^1$: it is continuous and piecewise $P_k$ on the Alfeld refinement, divergence-free, zero on $\partial T_F\supset F$, has $\int_F\dn v=0$, and $T_F\cap\Gh=F$ (Lemma~\ref{geom:one}), so no other face constraint or moment is affected. For unit $t\in T_{y_F}\Gamma$ apply Proposition~\ref{sh:joint} to $t':=P_Ft/|P_Ft|$, $P_F$ the projection onto the plane of $F$; $|P_Ft|\ge\frac12$ because the angle between $T_{y_F}\Gamma$ and $F$ is $O(\hG)$, and $t\cdot\M_F=|P_Ft|\,t'\cdot\M_F$ because $\M_F\parallel F$. Finally $|w^F|\ge c\ell^2|\hh_{y_F}|-C\ell^3$ (Lemma~\ref{geom:face}(d)).
\end{proof}
The rank computation was repeated independently, and the norm $\norm{\hat R}=91.678155469$ of the right inverse is certified (Appendix~\ref{app:cert}); for a regular macro-tetrahedron $\kappa_0=2.11\cdot10^{-3}$ (Proposition~\ref{ct:kappa}).

\begin{proposition}[(M3$^3$) on any mesh, $k\ge9$]\label{sh:bubble}
Let $K=T_F$, $\lambda_0$ its barycentric coordinate vanishing on $F$, $B:=(\lambda_1\lambda_2\lambda_3)^2$, $\tau\parallel F$ a unit vector and $w\ne0$ real weights. With $Q_w:=\sum_{i<j}w_{ij}\mu_i\mu_j$ written in $\lambda_1,\lambda_2,\lambda_3$, $\bar Q_w:=\int_FBQ_w/\int_FB$ and $\psi:=\lambda_0^2B(Q_w-\bar Q_w)(\nu\times\tau)$, the field $v:=\curl\psi\in P_9^3$, extended by zero, lies in $Y_h^1$ and
\[
\frac12\sum_{i<j}w_{ij}\int_F\mu_i\mu_j\partial_\nu v=-|\nabla\lambda_0|^2\,\tau\int_FB(Q_w-\bar Q_w)^2,\qquad\frac1{|F|}\int_FB(Q_w-\bar Q_w)^2\ge\frac{|w|^2}{8316000}.
\]
Hence (M3$^3$) holds for $k\ge9$ on every shape-regular mesh, with $\omega_F=T_F$ and $K_{\rm ov}=1$.
\end{proposition}

\begin{proof}
$\psi$ vanishes to second order on $\partial K$, so $v$ vanishes on $\partial K$, is continuous after extension by zero, and is divergence free. With $f:=\lambda_0^2B(Q_w-\bar Q_w)$, $\curl(f\Phi)=\nabla f\times\Phi+f\curl\Phi$, and on $F$ only the term with $2\lambda_0\nabla\lambda_0$ survives one normal derivative: $\partial_\nu v|_F=2|\nabla\lambda_0|^2B(Q_w-\bar Q_w)\,\nu\times(\nu\times\tau)=-2|\nabla\lambda_0|^2B(Q_w-\bar Q_w)\tau$, which has mean zero. The identity is $\int B(Q_w-\bar Q_w)Q_w=\int B(Q_w-\bar Q_w)^2$. The lower bound is $w^{\mathsf T}\mathcal Vw$, where $\mathcal V=\frac1{1188000}I-\frac1{2772000}(J-I)$ is the exact $B$-weighted covariance matrix of the three edge bubbles on $F$ ($J$ the all-ones matrix), with eigenvalues $1/831600$ (twice) and $1/8316000$. By scaling, $\norm{\nabla v}\le C|w|(\diam K)^{-1/2}$ and the moment is at least $c|w|^2|F|/h_K^2$, so the ratio is at least $c|w|(\diam F)^{1/2}$. Transfer to $t\in T_{y_F}\Gamma$ as in Theorem~\ref{sh:alfeld}.
\end{proof}

\begin{proposition}[$k=3$: one macro-element is not enough]\label{sh:k3single}
On the Alfeld split with $k=3$, $\dim\hat{\mathcal Y}_3=6$, $\dim\hat{\mathcal Y}^1_3=4$, and the image of $\hat E$ on $\hat{\mathcal Y}^1_3$ is the line spanned by $(\hat p_1-\hat p_2,\ \hat p_3-\hat p_1,\ \hat p_2-\hat p_3)$. Consequently, for a field supported in one macro-tetrahedron $T_F$, the physical moment $\M_F(v)$ is a multiple of $\sum_{\rm cyc}w_{ij}e_{ij}$; on the sphere this is a multiple of $R_{\pi/2}(x_c-o_F)$ (Lemma~\ref{geom:sag}(d)), which vanishes on equilateral faces, and it spans one direction only. So (M3$^3$) cannot be verified with $\omega_F=T_F$ for $k=3$. \status{computer-assisted, exact}
\end{proposition}

\subsection{Degree three}
For $k=3$ a single macro-element cannot work (Proposition~\ref{sh:k3single}): its fields see only one combination of the three edge weights, in one direction. The remedy is to use several macro-elements at once. Take a wall vertex $z$ of the face $F$ and the macro edge $[z,q_F]$ from $z$ to the apex $q_F$ of $T_F$; the \emph{edge star} $R(z,q_F)$ is the union of the macro-tetrahedra that contain this edge. Its wall faces all meet at $z$, and discrete divergence-free fields on the star can transport a gradient at $z$ between them. Appendix~\ref{app:k3} describes the traces of such fields exactly, computes the edge moments of an explicit family of ``vertex moves'' in closed form, and shows that, together with two further families of moves, they see every tangential direction, including at saddle and parabolic points of $\Gamma$. The result is the following theorem.

\begin{theorem}[(M3$^3$) for $k=3$; Theorem~\ref{k3:witness}]
On Alfeld refinements there are $\kappa_0=\kappa_0(\theta,\theta_0)>0$, $C_1$ and $h_\ast$ such that for $\hG\le h_\ast$, (M3$^3$) holds with $\omega_F$ an edge star at a vertex of $F$ and $K_{\rm ov},C_s\le C(\theta)$, where $\theta$ is the minimal angle of the macro mesh.
\end{theorem}
The constant $\kappa_0$ comes from compactness; on the sphere an explicit criterion is available (Theorem~\ref{k3:sphere}).

\begin{theorem}[Lower bound for $k\ge3$ on Alfeld refinements]\label{sh:k3main}
Under the assumptions of Section~\ref{sec:setting}, on an Alfeld refinement with $k\ge3$, if $|\hh|\,\partial_Nu\not\equiv0$ and $\hG\le h_1$, then for $\CNS$ ($\gamma\ge\gamma_2$) and for $\GS$ ($\gamma\ge\gamma_0$)
\[
\norm{\nabla(\ut-u_h)}_{\Oh}\ge c\,\hG^{3/2}\bigl\|\,|\hh|\,\partial_Nu\bigr\|_{L^2(\Gamma)}-C\hG^2,
\]
with $c=c(\theta,\theta_0,\Lambda_0)$. It covers saddle points, parabolic points and flat regions.
\end{theorem}

\begin{proof}
Theorem~\ref{k3:witness} and Theorem~\ref{sh:lb}; for $k>3$ use $Y_3\subset Y_k$.
\end{proof}
At a parabolic point the choice of vertex matters: for $\hh=(x+y)^2$ a fixed vertex can give a degenerate star (an exact computation; Appendix~\ref{app:repro}). This affects semidefinite regions such as cylinders.

\begin{corollary}[Sharpness of Scott's rate]\label{sh:sharp}
On Alfeld refinements with $k\ge3$, or on any mesh with $k\ge9$ satisfying (IS3), if $|\hh|\,\partial_Nu\not\equiv0$, $\gamma\in[\gamma_2,\gamma_{\max}]$ and $\epsd\le C\hG^{3/2}$, then for $\CNS$
\[
c\,\hG^{3/2}\le\norm{\nabla(\ut-u_h)}\le\tnorm{\ut-u_h}\le C\hG^{3/2}.
\]
The same holds for $\GS$ when $G=0$.
\end{corollary}

\begin{proof}
Theorems~\ref{er:D}, \ref{er:A} and~\ref{sh:lb} with Theorems~\ref{sh:alfeld}, \ref{sh:k3main} or Proposition~\ref{sh:bubble}.
\end{proof}

\begin{remark}[Worsey--Farin splits]\label{sh:rem:wf}
On Worsey--Farin splits each face of $\Gh$ is subdivided into three coplanar sub-faces of $\Th$ (so (M0$^3$) is used with this modification), and $\Yh^1$ then requires $\int\dn v=0$ on each sub-face; the moments $\M_F$ are still taken over the whole face $F$ of $\Gh$. The lower bound of Theorem~\ref{sh:lb} uses only the discrete momentum equation on $\Yh$, so it holds for every solution of $\GS$ or $\CNS$ that exists; but existence and the upper bounds are not available there, because (IS3) as formulated fails (Proposition~\ref{st:wf}). Proposition~\ref{sh:bubble} still applies for $k\ge9$, and Proposition~\ref{sh:wf} below gives single-macro witnesses for $k\ge4$ on barycentric Worsey--Farin splits (interior point at the barycentre, face points at the face barycentres; this split is affine-invariant, so the Piola arguments of Section~\ref{sh:sec:transport} apply, unlike the incentre-based split of \cite{GuzmanLischkeNeilan2022}). For $k=3$ a single Worsey--Farin macro-element has no witness, and Worsey--Farin edge stars give rank $2$ in all our tests (ranks modulo a prime; Appendix~\ref{app:repro}) \status{numerical}; the trace analysis of Appendix~\ref{sh:sec:k3} has not been redone for split faces.
\end{remark}

\begin{proposition}[Sharpness witnesses on barycentric Worsey--Farin splits]\label{sh:wf}
Let $\hat{\mathcal Y}_k$ be the continuous piecewise-$P_k$ divergence-free fields on the barycentric Worsey--Farin split of $\hat T$ vanishing on $\partial\hat T$, and $\hat{\mathcal Y}^1_k$ the subspace with zero mean of $\gh$ on each of the three sub-faces of $\hat F$.
\begin{enumerate}[label=(\roman*)]
\item $k=3$: $\dim\hat{\mathcal Y}_3=18$, $\dim\hat{\mathcal Y}^1_3=12$, and both the quadratic edge-moment map $\hat E$ and the first-moment map vanish identically on $\hat{\mathcal Y}^1_3$. There is no single-macro witness of degree $3$.
\item $k=4$: $\dim\hat{\mathcal Y}^1_4=72$, $\hat E$ is onto with $\norm{\hat R_{\rm WF}}=54.001475715$ and the first-moment map is onto with right-inverse norm $9.709257103$ (both certified). Hence (M3$^3$) holds on barycentric Worsey--Farin refinements for every $k\ge4$, with $\kappa_0$ as in Proposition~\ref{ct:kappa} and $\norm{\hat R}$ replaced by $54.0015$, and Theorem~\ref{sh:lb} gives the lower bound $c\,\hG^{3/2}\norm{|\hh|\partial_Nu}_{L^2(\Gamma)}-C\hG^2$ for every solution of $\GS$ or $\CNS$ that exists.
\end{enumerate}
\status{computer-assisted, exact}
\end{proposition}
For $k\ge4$ this is a sharpness witness: together with an upper bound it would give the exact rate $\hG^{3/2}$ on barycentric Worsey--Farin splits. The upper bound needs the reformulated stability input of Remark~\ref{st:rem:status}, which we do not have.
